\section{Proposed Approach (RQ1)}
\label{sec:approach}

\begin{figure*}[t]
    \centering
    \includegraphics[width=\textwidth]{SchémaFinal.png}
    \caption{Overview of the proposed approach}
    \label{fig:approach}
\end{figure*}

To address the difficulty of identifying business-rule implementations within large and complex codebases, we introduce an automated approach designed to assist developers throughout this process. This section provides an overview of the proposed approach and details the successive steps involved in its execution.

\subsection{Overview of the Approach}

To address the challenge of locating business rules in large codebases, we propose an approach that progressively narrows the search space by exploiting the business concepts expressed in rule descriptions. Specifically, business rule descriptions contain business terms that refer to core business concepts, which are represented in the source code by business entities. In this work, we consider business entities to include business classes and their attributes. The objective is therefore to identify terms in a business rule description that can serve as links to the source code. When these terms correspond to business entities that represent the concepts described in the rule, they provide a basis for connecting the textual description of the rule to its implementation in the source code.

To establish this link between the business rule description and its implementation in the source code, our approach combines information extracted from databases with information derived from source code to automatically identify relevant business classes. It then associates the business terms appearing in a given rule description with the corresponding business entities and uses these associations to identify methods that manipulate the relevant entities. The result is a set of candidate methods that are likely to participate in the implementation of the rule. Our goal is not to identify a single implementation point, but rather to provide developers with a focused set of code elements for further investigation.

The approach consists of two preliminary steps, performed once for the entire codebase, and three rule-specific steps, applied for each business rule. All these steps are automated.

\begin{itemize}
    \item \textbf{Preliminary Steps}
    \begin{enumerate}
        \item[\textbf{P.1}] Identify business entities in the source code.
        \item[\textbf{P.2}] For each business entity, identify the methods that manipulate it.
    \end{enumerate}

    \item \textbf{Rule-Specific Steps}
    \begin{enumerate}
        \item[\textbf{R.1}] Extract business vocabulary from the business rule description.
        \item[\textbf{R.2}] Map business vocabulary to business entities.
        \item[\textbf{R.3}] Select P.2 methods using business entities identified in R.2.
    \end{enumerate}
\end{itemize}

The objective of these steps is to obtain a set of candidate methods that manipulate the business entities identified in R.2 and may therefore contribute to the implementation of the given business rule. Since business entities are identified from both classes and attributes in the source code, different entities may have closely related or even similar names. As a result, it may not be possible to determine with certainty which entity best represents a given business term. Therefore, for each business term, we consider all business entities that may potentially represent the corresponding business concept. In R.3, we use an OR condition when filtering methods based on these entities, retaining any method that manipulates at least one of them. This inclusive strategy increases the size of the candidate set, but reduces the risk of excluding methods that may contribute to the implementation of the business rule and thus improves the likelihood of covering all relevant implementation methods.

\subsection{Detailed Steps of the Approach}

The remainder of this section describes each step in detail and explains how it contributes to the identification of methods that potentially implement a given business rule. Figure~\ref{fig:approach} provides an overview of the proposed approach and its main steps.

\vspace{-0.3em}

\subsubsection{Preliminary Steps}\noindent

\vspace{-0.3em}

\paragraph{P.1\quad Identify business entities in the source code}\noindent

We begin by identifying the business entities present in the source code, starting with business classes. Our goal is to identify the classes that represent core business concepts and manage data relevant to the application domain. In our approach, these classes are represented by classes belonging to the Object-Relational Mapping (ORM) layer, which provide access to data persisted in the database. As discussed in Section~\ref{sec:businessVar}, business variables also provide useful information for locating business concepts in source code. However, directly identifying such variables in a large industrial system would require examining a potentially large number of program elements. We therefore use business classes as the starting point for identifying business entities, since they provide a more stable and structured representation of the domain concepts implemented in the system.

To automate their identification, we rely on information available from both the source code and the underlying database. More specifically, we consider as business classes the ORM classes that expose or manipulate data persisted in the database. This criterion allows us to identify the classes that represent the application's persistent business entities without requiring manual inspection. For each identified business class, we also consider its attributes, as they represent the data elements through which business concepts are concretely represented in the source code. 

At the end of this step, we obtain a set of business entities composed of the identified business classes and their corresponding attributes.

\paragraph{P.2\quad For each business entity, identify the methods that manipulate it}\noindent

In this step, we focus on identifying, for each business entity, the methods that manipulate or reference it. We consider two types of business entities: classes and attributes.

For each business class, we collect the methods that reference the class, as these methods may implement operations involving the corresponding business entity. 

For attributes, since they are private, access to them is typically performed through their accessors, namely getter and setter methods. Therefore, for each business attribute, we identify the methods that invoke its getter or setter methods.

At the end of this step, we obtain a collection that associates each business entity with a list of methods.

\vspace{-0.5em}

\subsubsection{Rule-Specific Steps}\noindent

\vspace{-0.5em}

\paragraph{R.1\quad Extract business vocabulary from the business rule description}\noindent

Given a description of a business rule, we aim to identify the business vocabulary it refers to. In a context where no dictionary is available to define in advance a list of business vocabulary specific to the domain, identifying business terms directly from the rule description is particularly challenging. Indeed, there is no clear boundary between domain-specific business terms and general words, and the same term may have different meanings depending on the context in which it appears. Moreover, a business rule may contain various types of words, including domain concepts, actions, conditions, and contextual information, which makes it difficult to determine which terms are relevant for locating their implementation. Therefore, the identification process must distinguish terms that convey business concepts from words that do not provide useful information.

Moreover, NLP-based approaches cannot reliably identify all relevant business terms because they do not necessarily belong to a specific linguistic category. Business terms may correspond to individual nouns, noun phrases, or other linguistic structures, making it difficult to define a single set of linguistic rules for their extraction.

Automatically identifying business terms from the rule description may lead to two opposite problems: retaining words that are not relevant to the business domain or discarding terms that are actually useful for locating the rule implementation. We therefore adopt a conservative strategy. Rather than selecting only the terms that are initially considered to be business terms, we retain a broad set of possible candidates extracted from the rule description.

For each business rule description, we first remove stop words and then generate chunks corresponding to consecutive sequences of the remaining words. These chunks cover different combinations of words that may form a business term, ranging from individual words to longer expressions. This strategy allows us to preserve potential business terms even when their linguistic structure cannot be determined in advance. As a result, the generated list may contain both relevant chunks that represent business concepts and irrelevant ones that do not correspond to meaningful business terms. The latter are filtered out in the subsequent steps using additional information available from the source code. At this stage, the objective is therefore not to precisely identify business terms, but rather to preserve the information contained in the rule description and avoid prematurely excluding potentially relevant candidates.

For example, consider the rule description \textit{``If specified, the issuing budget must exist for the local authority.''} After removing stop words, the remaining words are \textit{specified}, \textit{issuing}, \textit{budget}, \textit{exist}, \textit{local}, and \textit{authority}. We then generate consecutive chunks from these words, such as \textit{budget}, \textit{issuing budget}, \textit{local}, and \textit{local authority}. At this stage, all generated chunks are retained as candidates, without determining whether they correspond to actual business terms. The filtering of irrelevant candidates is performed in the subsequent steps.

At the end of this step, all generated chunks are retained as candidates, regardless of the amount or relevance of information they contain. The main objective at this stage is to preserve all potentially useful information from the rule description and ensure that no relevant information is discarded before the subsequent steps.

\vspace{-0.5em}

\paragraph{R.2\quad Map business vocabulary to business entities}\noindent

\vspace{-0.1em}

Each chunk obtained from the previous step, P.1, is represented as an embedding to capture its semantic and linguistic information. Similarly, all business entities identified in P.1 are represented by embeddings of their names. For each chunk, we compare its embedding with the embeddings of all business entities by computing the cosine similarity between the corresponding vectors. The resulting similarity scores are then used to rank the business entities according to their semantic similarity to the chunk. Thus, each chunk is associated with a ranked list of business entities, from the most similar to the least similar.

We then apply two parameters to filter the chunks and their associated business entities. First, we use a similarity threshold to determine whether a chunk is sufficiently related to at least one business entity. The cosine similarity score ranges from 0 to 1, with higher values indicating greater semantic similarity between a chunk and a business entity. We discard a chunk if its highest similarity score is below the selected threshold, as this indicates that the chunk is not sufficiently related to any business entity. After experimenting with different threshold values, we selected a threshold of 0.7.

Second, for each remaining chunk, we determine how many similarity scores should be retained. We experimented with retaining the entities corresponding to the top-1, top-2, and top-3 similarity scores. Based on these experiments, we selected the top-3 similarity scores. Thus, for each chunk, we retain the business entities associated with the three highest similarity scores. This choice allows us to preserve several potentially relevant business entities while limiting the amount of irrelevant information introduced at this stage.

At the end of this step, each retained chunk is associated with the business entities corresponding to its three highest similarity scores.

\vspace{-0.5em}

\paragraph{R.3\quad Select P.2 methods using business entities identified in R.2}\noindent

\vspace{-0.1em}

The goal of this step is to filter the methods in the source code and retain only those that are likely to implement a given business rule. To this end, we rely on the results of Step (P.2), which provides a list of methods associated with each business entity, as well as the list of business entities identified from the rule description in Step (R.2).

We retain methods associated with at least one business entity from the list obtained from the chunks of the rule description. This process removes methods with no apparent connection to the rule ( uses other business entities ) , resulting in a smaller set of candidate methods for further analysis. Consequently, developers do not need to inspect the entire codebase but can focus their investigation on methods that are potentially related to the implementation of the business rule.

The primary objective of this step is to maintain a high recall rather than to maximize precision. Since methods discarded at this stage cannot be recovered in others steps, it is essential to ensure that the method implementing the business rule remains among the retained candidates. Some irrelevant methods may therefore be preserved, even if they do not ultimately contribute to the implementation of the rule. This trade-off is intentional: the filtering process should reduce the size of the search space while minimizing the risk of excluding the actual implementation method. Thus, even when the number of candidate methods is substantially reduced, we aim to ensure that the method implementing the business rule remains within the resulting set of candidates.